\chapter{Osterwalder--Schrader Reconstruction}
\label{ch:volume-iv-osterwalder-schrader-reconstruction}
\ConventionsMixed

\label{ch:os-reconstruction}

\section{Euclidean Correlation Data}

Let \(E=\mathbb R^D\) be Euclidean space with coordinates
\[
  x=(\tau,\mathbf x),\qquad \mathbf x\in\mathbb R^{D-1}.
\]
Let \(V\) be the complex vector space of field labels.  For scalar fields
\(V=\mathbb C\).  For tensor, spinor, or internal multiplets, \(V\) carries the
corresponding Euclidean rotation and internal representations.  An \(n\)-point
Euclidean test insertion is an element of
\[
  \mathcal S(E^n;V^{\otimes n}),
\]
or of \(C_c^\infty(E^n;V^{\otimes n})\) if compactly supported test functions
are the chosen test class.  The test class is part of the data because it
determines the topology in which the Schwinger functions are distributions.
Let
\[
  \Delta_n
  =
  \bigcup_{1\le i<j\le n}\{(x_1,\ldots,x_n)\in E^n\mid x_i=x_j\}
\]
be the union of partial diagonals.  The corrected OS reconstruction theorem is
most honestly stated on the zero-diagonal Schwartz space
\[
  \mathcal S_\circ(E^n;V^{\otimes n})
  =
  \{f\in\mathcal S(E^n;V^{\otimes n})
    \mid \partial^\alpha f|_{\Delta_n}=0\text{ for every multi-index }\alpha\}.
\]
Pairing with \(\mathcal S_\circ\) treats the Schwinger functions as
distributions away from coincident points.  A choice of contact-term extension
to the full Schwartz space is extra structure and is not part of the OS
reconstruction input used here.

\begin{definition}[Schwinger hierarchy]
A Schwinger hierarchy with field-label space \(V\) is a family
\[
  S_n\in \mathcal S_\circ'(E^n;(V^*)^{\otimes n}),
  \qquad n\ge0,
\]
with \(S_0=1\).  The value \(S_n(f_n)\) denotes the distributional pairing of
\(S_n\) with the \(V^{\otimes n}\)-valued test function \(f_n\).  In
component notation this is written as
\[
  S_{\alpha_1\cdots\alpha_n}(x_1,\ldots,x_n),
  \qquad \alpha_j\in I,
\]
where \(I\) is a basis of labels for \(V\).  The component notation is a
kernel notation for the distributional pairing; the pairing with test
functions is the primary object.
\end{definition}

In a construction by a Euclidean measure, these distributions are moments of
random distributions.  In an axiomatic Euclidean formulation, the hierarchy
\(\{S_n\}_{n\ge0}\) is primary data.

Let \(\theta\) be time reflection:
\[
  \theta(\tau,\mathbf x)=(-\tau,\mathbf x).
\]
The Euclidean group is
\[
  \operatorname{Euc}(D)=\mathbb R^D\rtimes O(D).
\]
It acts on test insertions by pullback on \(E^n\) and by the declared label
representation on \(V\).
The reflection also acts on labels.  We denote by \(J:V\to V\) the antilinear
label map induced by field adjunction and by the Euclidean reflection of
tensor or spinor indices.  For a scalar Hermitian field, \(J\) is ordinary
complex conjugation.  For a test function \(f_n\) supported in the ordered
positive-time region
\[
  E^n_>=\{(x_1,\ldots,x_n):0<\tau_1<\cdots<\tau_n\},
\]
define the reflected adjoint test function \(\Theta f_n\) by reversing the
order, reflecting all arguments, and applying \(J\) to the labels.  For a
scalar test function this reads
\[
  (\Theta f_n)(x_1,\ldots,x_n)
  =
  \overline{f_n(\theta x_n,\ldots,\theta x_1)}.
\]
With labels, the same formula is composed with \(J^{\otimes n}\).  Let
\(\Theta\) denote the same antilinear operation on arbitrary \(n\)-point test
functions; positive-time support is needed only when we form the OS
pre-Hilbert space below.  Let
\(\mathcal E_+\) be the vector space of finite sequences
\[
  F=(f_0,f_1,f_2,\ldots),
\]
where only finitely many \(f_n\) are nonzero and every \(f_n\) is a labelled
zero-diagonal test function supported in positive Euclidean time with a chosen
ordering of its time variables.  Reflection positivity is the requirement that the
sesquilinear form
\begin{equation}
  \langle F,G\rangle_{\rm OS}
  =
  \sum_{m,n} S_{m+n}\bigl(\Theta f_m,g_n\bigr)
  \label{eq:os-positive-time-form}
\end{equation}
satisfy
\begin{equation}
  \langle F,F\rangle_{\rm OS}\ge0
  \label{eq:os-reflection-positivity}
\end{equation}
for every \(F\in\mathcal E_+\).

The notation \(S_{m+n}(\Theta f_m,g_n)\) means that the reflected \(m\)
negative-time insertions are followed by the \(n\) positive-time insertions in
the argument of the distribution.  Thus the object required to be positive is
the reflected positive-time sesquilinear form on \(\mathcal E_+\).

\begin{definition}[OS-admissible Euclidean data]
\label{def:os-admissible-euclidean-data}
An OS-admissible Euclidean field presentation is a Schwinger hierarchy
\(\{S_n\}_{n\ge0}\), together with the label reflection \(J\), satisfying:
\begin{enumerate}[label=(\roman*)]
  \item normalization: \(S_0=1\);
  \item Euclidean covariance: \(S_n(gf_n)=S_n(f_n)\) for
  \(g\in \operatorname{Euc}(D)\);
  \item permutation symmetry: \(S_n\) is invariant under simultaneous
  permutation of arguments and labels;
  \item Hermiticity: for every \(n\ge0\) and every labelled \(n\)-point test
  function \(f_n\),
  \begin{equation}
    S_n(\Theta f_n)
    =
    \overline{S_n(f_n)} .
    \label{eq:os-hermiticity}
  \end{equation}
  \item reflection positivity:
  \(\langle F,F\rangle_{\rm OS}\ge0\) for all \(F\in\mathcal E_+\);
  \item OS-II regularity and linear growth: the hierarchy has the
  positive-time semigroup continuity, contractivity, and linear-growth bounds
  stated below.
\end{enumerate}
For fermionic or graded fields, the permutation line is replaced by graded
permutation symmetry with the Koszul sign determined by the declared parity of
each label.  Clustering is added when the target Lorentzian theory is required
to have a unique vacuum vector.
\end{definition}

The last line is a load-bearing analytic hypothesis, not a cosmetic
regularity clause.  Let \(S_n^{\rm red}\) denote the translation-reduced
Schwinger distribution in \(n-1\) difference variables, acting on the
corresponding zero-diagonal difference-test space.  With a fixed family of
Schwartz seminorms
\[
  q_{n,N}(F)
  =
  \sum_{|\alpha|,|\beta|\le N}
  \sup_{\Xi\in E^{n-1}}
  \left|\Xi^\beta\partial^\alpha F(\Xi)\right|,
  \qquad F\in\mathcal S_\circ(E^{n-1};V^{\otimes n}),
\]
one convenient OS-II linear-growth condition is the existence of constants
\[
  C_{\rm OS},N_0\in\mathbb N,\qquad
  A,B,\gamma>0,
\]
independent of \(n\), such that
\begin{equation}
  |S_n^{\rm red}(F)|
  \le
  A\,B^n\,(n!)^\gamma\,
  q_{n,N_0+C_{\rm OS}n}(F)
  \label{eq:os-linear-growth-condition}
\end{equation}
for all \(F\) in that test space.  Equivalent OS-II formulations package the
same content by using a fixed seminorm family \(|\cdot|_{s,n}\) whose
definition already contains the \(n\)-dependence, together with constants
\(\sigma_n\le A B^n(n!)^\gamma\).  The word ``linear'' refers to the allowed
growth of the distributional seminorm order with the number of insertions.  It
is not a statement that a fixed \(n\)-point function grows linearly in
spacetime.

Together with the positive-time semigroup hypotheses
\eqref{eq:os-time-semigroup-law}--\eqref{eq:os-time-strong-continuity}, this
condition supplies three analytic facts used later: positive-time translations
act continuously on the quotient, Euclidean rotations can be analytically
continued to Lorentz boosts, and the continued correlation functions have
boundary values in the declared target distribution class, for example
tempered distributions when the Schwinger functions obey
\eqref{eq:os-linear-growth-condition}.  The correction from the first
Osterwalder--Schrader reconstruction paper to the second is precisely that
reflection positivity, Euclidean covariance, symmetry, clustering, and
ordinary temperedness do not by themselves supply this continuation theorem.
Thus an application of OS reconstruction must either prove an OS-II growth
condition such as \eqref{eq:os-linear-growth-condition}, or replace it by a
different theorem that supplies the same Lorentzian boundary-value estimates.
\footnote{See K.~Osterwalder and R.~Schrader, ``Axioms for Euclidean Green's
Functions,'' Commun. Math. Phys. \textbf{31} (1973) 83--112, and ``Axioms for
Euclidean Green's Functions II,'' Commun. Math. Phys. \textbf{42} (1975)
281--305.  The second paper states the corrected reconstruction theorem with
the linear-growth input.  Modern checks of this analytic input are model
dependent.  For example, Gubinelli and Hofmanov\'a, ``A PDE Construction of
the Euclidean \(\Phi^4_3\) Quantum Field Theory,'' Commun. Math. Phys.
\textbf{384} (2021) 1--75, construct limit measures with reflection
positivity, translation invariance, non-Gaussianity, and stretched exponential
integrability estimates strong enough for the OS distributional regularity
axiom in their formulation; rotation invariance and clustering remain separate
inputs in that construction.}

\begin{remark}[CFT reconstruction is a different route]
For a Euclidean CFT, one should not silently replace the OS-II growth package
by conformal terminology.  Conformal kinematics gives additional structure:
radial ordering, the inversion pairing, positivity of the state-operator
inner product, and conformal covariance of operator insertions.  In that
setting a Lorentzian theory can be reconstructed by proving the convergence
and distributional boundary-value properties of conformal correlation
functions directly from the conformal Hilbert-space and OPE data.  This is a
different theorem from OS reconstruction.  Kravchuk, Qiao, and Rychkov prove
the relevant Lorentzian distribution and OPE-convergence estimates for CFT
correlators using conformal kinematics and radial variables;\footnote{See
P.~Kravchuk, J.~Qiao, and S.~Rychkov, ``Distributions in CFT II. Minkowski
Space,'' JHEP \textbf{08} (2021) 094, arXiv:2104.02090.} the CFT volume uses
that route only after the state-operator Hilbert space and conformal
representation theory have been stated as hypotheses.
\end{remark}

\begin{figure}[H]
  \centering
  \resizebox{0.82\textwidth}{!}{%
  \begin{tikzpicture}[x=1cm,y=1cm,every node/.style={font=\scriptsize}]
    \tikzset{
      axis/.style={-{Latex[length=2mm]},line width=0.75pt},
      dot/.style={circle,draw,fill=blue!12,line width=0.8pt,minimum size=0.34cm},
      rdot/.style={circle,draw,fill=orange!14,line width=0.8pt,minimum size=0.34cm},
      arrow/.style={-{Latex[length=2mm]},line width=0.85pt}
    }
    \draw[axis] (-4.2,0)--(4.2,0) node[right] {\(\tau\)};
    \draw[axis] (0,-2.0)--(0,2.0) node[above] {space};
    \draw[line width=0.8pt,dashed] (0,-1.75)--(0,1.75);
    \node[dot] (a) at (1.15,0.95) {};
    \node[dot] (b) at (2.55,-0.55) {};
    \node[rdot] (ta) at (-1.15,0.95) {};
    \node[rdot] (tb) at (-2.55,-0.55) {};
    \draw[arrow,orange!80!black] (a)--(ta)
      node[midway,above] {\(\theta\)};
    \draw[arrow,orange!80!black] (b)--(tb)
      node[midway,below] {\(\theta\)};
    \node[blue!55!black] at (2.15,1.55) {positive-time tests};
    \node[orange!80!black] at (-2.25,1.55) {reflected tests};
    \node at (0,-2.25)
      {\(\sum S(\Theta f,f)\ge0\) becomes the Lorentzian inner product.};
  \end{tikzpicture}
  }
  \caption{Reflection positivity turns Euclidean reflection through
  \(\tau=0\) into Hilbert-space positivity after reconstruction.}
  \label{fig:os-reflection-positivity}
\end{figure}

\section{Hilbert Space From Positive Time}

Use the positive-time vector space \(\mathcal E_+\) and the form
\eqref{eq:os-positive-time-form}.  Reflection positivity says that this form
is nonnegative.  Let
\[
  \mathcal N=\{F\in\mathcal E_+\mid \langle F,F\rangle_{\rm OS}=0\}.
\]
The reconstructed Hilbert space is
\[
  \Hilb_{\rm OS}
  =
  \overline{\mathcal E_+/\mathcal N}.
\]
The vector represented by \(F=(1,0,0,\ldots)\) is the vacuum.
The quotient is taken before completion.  Since
\eqref{eq:os-positive-time-form} is positive semidefinite, the
Cauchy--Schwarz inequality holds on \(\mathcal E_+\), and
\(\mathcal N\) is exactly the subspace of zero-norm vectors.  Therefore
\(\mathcal E_+/\mathcal N\) is a pre-Hilbert space with inner product induced
by \(\langle\cdot,\cdot\rangle_{\rm OS}\).

Euclidean time translations by \(s\ge0\) preserve positive-time support and
therefore act on \(\mathcal E_+\).  Write \(a_s(\tau,\mathbf x)=(\tau+s,\mathbf
x)\), and let \(R_s\) be the induced action on test sequences, so that each
component is translated by \(a_s\).  The OS regularity hypotheses needed at
this point are the following seminorm statements on \(\mathcal E_+\):
\begin{align}
  R_0&=\mathbf 1,
  &
  R_{s+t}&=R_sR_t,
  \label{eq:os-time-semigroup-law}\\
  \|R_sF\|_{\rm OS}&\le \|F\|_{\rm OS},
  \label{eq:os-time-contractivity}\\
  \langle R_sF,G\rangle_{\rm OS}
  &=
  \langle F,R_sG\rangle_{\rm OS},
  \label{eq:os-time-symmetry}\\
  \lim_{s\downarrow0}\|R_sF-F\|_{\rm OS}&=0,
  \label{eq:os-time-strong-continuity}
\end{align}
for all \(F,G\in\mathcal E_+\), where
\(\|F\|_{\rm OS}^2=\langle F,F\rangle_{\rm OS}\).  The inequality
\eqref{eq:os-time-contractivity} is the OS-II contractivity strengthening.  It
is an analytic hypothesis on the reflection-positive seminorm, additional to
the algebraic preservation of positive-time support and to Euclidean covariance
of the Schwinger functions.  This is the point at which the Klein--Landau
regularity gap in the original reconstruction argument is closed.

\begin{remark}[Failure mode of the first OS reconstruction theorem]
\label{rem:klein-landau-gap}
The OS-I axioms do not by themselves imply that \(R_s\) is bounded on the
reflection-positive quotient.  The obstruction is concrete: even when the
Euclidean hierarchy is reflection positive and Euclidean covariant, the
seminorm of \(R_sF\) may fail to be controlled by the seminorm of \(F\), so a
zero-norm vector can be moved to a nonzero class or a Cauchy sequence can be
moved to a non-Cauchy sequence.  In that situation the formula
\([F]\mapsto[R_sF]\) does not define a contraction on \(\Hilb_{\rm OS}\), and
there is no self-adjoint Hamiltonian obtained from positive Euclidean time
translations.  The added OS-II contractivity and continuity assumptions
\eqref{eq:os-time-contractivity}--\eqref{eq:os-time-strong-continuity} are
exactly the hypotheses that remove this failure mode.
\end{remark}

Contractivity implies \(R_s\mathcal N\subset\mathcal N\), so \(R_s\) descends to
a bounded operator \(T_E(s)\) on the quotient and then to \(\Hilb_{\rm OS}\).
Equations
\eqref{eq:os-time-semigroup-law}--\eqref{eq:os-time-strong-continuity} say
precisely that \(\{T_E(s)\}_{s\ge0}\) is a strongly continuous semigroup of
self-adjoint contractions.
Indeed, \eqref{eq:os-time-symmetry} gives, on the dense quotient,
\[
  \langle [F],T_E(s)[G]\rangle_{\rm OS}
  =
  \langle [F],[R_sG]\rangle_{\rm OS}
  =
  \langle [R_sF],[G]\rangle_{\rm OS}
  =
  \langle T_E(s)[F],[G]\rangle_{\rm OS}.
\]
The equality descends to the quotient because \(R_s\mathcal N\subset
\mathcal N\), and it extends to \(\Hilb_{\rm OS}\) because \(T_E(s)\) is
bounded by \eqref{eq:os-time-contractivity}.  Thus the self-adjointness of
the Euclidean time semigroup is not inferred from a formal reflection
manipulation outside positive time; it is precisely the OS-II symmetry
condition on the reflection-positive quotient.

Let \(A\) be its infinitesimal generator,
\[
  A\psi=\lim_{s\downarrow0}\frac{T_E(s)\psi-\psi}{s},
\]
on the vectors for which the limit exists in \(\Hilb_{\rm OS}\).  The
self-adjoint contraction-semigroup theorem gives \(A=A^*\) and
\(\sigma(A)\subset(-\infty,0]\).  Setting \(H=-A\), the spectral theorem gives
\[
  T_E(s)=\ee^{-sH},\qquad H=H^*,\qquad H\ge0 .
\]
Stone's theorem then applies to the strongly continuous unitary group
\(U_{\rm time}(t)=\ee^{-\ii tH}\): its named hypotheses are the unitary group
law, strong continuity in \(t\), and self-adjointness of the generator.  This
\(H\) is the Lorentzian Hamiltonian in the reconstructed vacuum sector.
Spatial translations and rotations act unitarily on the quotient.
Together with the analytic continuation of Euclidean rotations, the full OS
theorem gives a strongly continuous unitary representation of the connected
Poincare group with joint translation spectrum contained in the closed forward
cone.

The passage from Euclidean covariance to Lorentzian covariance occurs here:
the Lorentzian Hilbert space and Poincare representation are constructed from
the positive-time quotient and the analytic-continuation data.

\begin{figure}[H]
  \centering
  \resizebox{0.94\textwidth}{!}{%
  \begin{tikzpicture}[x=1cm,y=1cm,every node/.style={font=\scriptsize}]
    \tikzset{
      box/.style={draw,rounded corners=4pt,line width=0.8pt,fill=blue!4,
        minimum width=3.25cm,minimum height=0.90cm,align=center},
      arrow/.style={-{Latex[length=2mm]},line width=0.85pt}
    }
    \node[box] (sch) at (-5.25,0.45) {Schwinger hierarchy\\\(\{S_n\}\)};
    \node[box] (pos) at (-1.75,0.45) {positive-time\\test sequences};
    \node[box] (hilb) at (1.75,0.45) {OS quotient\\\(\Hilb_{\rm OS}\)};
    \node[box] (lor) at (5.25,0.45)
      {Wightman data\\\(\Hilb_{\rm OS},\vac_{\rm OS},U,\widehat\Phi_\alpha\)};
    \node[box,fill=green!7] (semigroup) at (0,-1.45)
      {Euclidean time semigroup\\\(\ee^{-sH}\)};
    \draw[arrow] (sch)--(pos);
    \draw[arrow] (pos)--(hilb);
    \draw[arrow] (hilb)--(lor);
    \draw[arrow] (pos)--(semigroup);
    \draw[arrow] (semigroup)--(hilb);
    \node[blue!60!black] at (-3.50,1.32) {reflection positivity};
    \node[blue!60!black] at (3.50,1.32) {analytic continuation};
  \end{tikzpicture}
  }
  \caption{OS reconstruction turns Euclidean correlation data into a
  Lorentzian Hilbert-space theory through a positive-time quotient.}
  \label{fig:os-reconstruction-flow}
\end{figure}

\section{Fields on the Reconstructed Domain}

Let \(\mathcal D_{\rm OS}\subset\Hilb_{\rm OS}\) be the image of
\(\mathcal E_+\) in the quotient.  This dense subspace is the natural common
domain for the reconstructed fields before closure.  If \(h\) is a Euclidean
test function valued in \(V\), with support at positive time and ordered
before the supports appearing in \(F\in\mathcal E_+\), define
\[
  \Phi_E(h)[F]
  =
  [h\otimes F],
\]
where \(h\otimes F\) denotes the sequence obtained by inserting one additional
Euclidean field smeared with \(h\).  Permutation symmetry of the Schwinger
functions and reflection positivity make this definition compatible with the
null quotient whenever the supports obey the stated ordering.  Under the
operator-domain hypotheses included in the reconstruction theorem, limits of
ordered supports give closable Euclidean field operators on
\(\mathcal D_{\rm OS}\).

\begin{proposition}[Closability of ordered Euclidean insertions]
\label{prop:os-ordered-field-closability}
Assume the operator-domain part of the OS data supplies the following
adjoint-insertion identity.  For every ordered positive-time test insertion
\(h\) there is an ordered test insertion \(h^\sharp\), with adjoint label and
with supports arranged by a common positive time translation, such that for
all \(F,G\in\mathcal E_+\)
\begin{equation}
  \langle [G],\Phi_E(h)[F]\rangle_{\rm OS}
  =
  \langle \Phi_E(h^\sharp)[G],[F]\rangle_{\rm OS}.
  \label{eq:os-ordered-insertion-adjoint-identity}
\end{equation}
Then \(\Phi_E(h)\) is closable on the dense domain
\(\mathcal D_{\rm OS}\).
\end{proposition}

\begin{proof}
Let \(\psi_k\in\mathcal D_{\rm OS}\) satisfy
\(\psi_k\to0\) and \(\Phi_E(h)\psi_k\to\eta\) in \(\Hilb_{\rm OS}\).  To
prove closability it suffices to show \(\eta=0\).  For a dense test vector
\([G]\in\mathcal D_{\rm OS}\), the adjoint-insertion identity gives
\[
  \langle [G],\eta\rangle_{\rm OS}
  =
  \lim_k\langle [G],\Phi_E(h)\psi_k\rangle_{\rm OS}
  =
  \lim_k\langle \Phi_E(h^\sharp)[G],\psi_k\rangle_{\rm OS}
  =
  0,
\]
because \(\Phi_E(h^\sharp)[G]\) is a fixed vector in
\(\mathcal D_{\rm OS}\) and \(\psi_k\to0\).  Since \(\mathcal D_{\rm OS}\) is
dense, \(\eta\) is orthogonal to all of \(\Hilb_{\rm OS}\), hence
\(\eta=0\).  This is the standard closability criterion, with the
QFT-specific content isolated in the matrix-element identity
\eqref{eq:os-ordered-insertion-adjoint-identity}.
\end{proof}

Lorentzian fields are obtained by continuing these matrix elements on the
reconstructed Hilbert space.  For ordered imaginary parts
\[
  \epsilon_1>\epsilon_2>\cdots>\epsilon_n>0,
\]
the Wightman distributions are boundary values of the Schwinger functions
continued to
\begin{equation}
  \tau_j=\epsilon_j+\ii t_j,
  \qquad j=1,\ldots,n.
  \label{eq:os-wightman-boundary-values}
\end{equation}
The reconstructed operator \(\widehat\Phi_\alpha(t,\mathbf f)\) is the
operator-valued distribution whose vacuum matrix elements are these boundary
values and whose time dependence is implemented by the unitary group
\[
  \ee^{-\ii tH}.
\]
Thus the OS construction produces the vacuum vector, the Hilbert space, the
positive Hamiltonian, and the field domains from the same Euclidean
correlation hierarchy.

As in the Wightman presentation, these reconstructed fields are generally
unbounded distributional coordinates.  A represented local observable algebra
\(\mathcal R(\mathcal O)\) is obtained by specifying a closure, affiliation,
or bounded functional-calculus construction, or by building a local net
directly and proving that the reconstructed fields are affiliated with it.
The OS theorem reconstructs a Lorentzian Wightman field
presentation; comparison with AQFT local observable content is a further
construction with its own hypotheses.

\section{Analytic Continuation Input}

The OS-II regularity assumptions are used through the following
boundary-value package.  It is stated separately because this is the part of
the reconstruction theorem where the first OS paper needed correction.

\begin{definition}[OS Lorentzian boundary-value package]
\label{def:os-lorentzian-boundary-value-package}
For each \(n\), the Schwinger distribution \(S_n\) is said to satisfy the OS
Lorentzian boundary-value package if the following data and estimates are
available.

\begin{enumerate}[label=(\roman*)]
\item For every Euclidean time ordering there is a holomorphic function
\[
  \mathcal W_n(z_1,\ldots,z_n)
\]
on the corresponding forward tube in the difference variables
\(\operatorname{Im}(z_j-z_{j+1})\in V_+^\circ\), where \(V_+^\circ\) is the
open Lorentzian forward cone, whose restriction to ordered imaginary times
agrees with \(S_n\) away from coincident points.
\item The holomorphic functions obey tempered tube bounds: for every compact
subcone \(K\subset (V_+^\circ)^{n-1}\) and every positive lower bound on the
distance of the imaginary parts from the boundary of the cone, the boundary
pairings with Schwartz tests are bounded by finitely many Schwartz seminorms
with polynomial dependence on the reciprocal distance to the cone boundary.
\item The permuted Euclidean orderings give holomorphic functions on the
corresponding permuted tubes that agree on their common Euclidean real
submanifolds.  For spacelike separated adjacent insertions, the two
permuted tubes have a connected edge-of-the-wedge continuation through the
corresponding Jost domain, with the Koszul sign included for graded fields.
\end{enumerate}
\end{definition}

The linear-growth estimate \eqref{eq:os-linear-growth-condition} is one
standard route to this package in the OS-II theorem: the controlled growth of
distributional seminorm order is what keeps the Fourier--Laplace continuation
inside the tempered distribution class as the imaginary parts approach the
Lorentzian boundary.  In a concrete constructive model, this route is a
separate analytic estimate to be proved, not a formal consequence of
reflection positivity alone.

\begin{proposition}[Consequences of the OS boundary-value package]
\label{prop:os-boundary-package-consequences}
Assume Definition~\ref{def:os-lorentzian-boundary-value-package}.  Then:
\begin{enumerate}[label=(\alph*)]
\item The limits of \(\mathcal W_n(x+\ii y)\) as
\(y\to0\) within the ordered tube exist in
\(\mathcal S'((\mathbb M^D)^n)\) and define tempered Wightman distributions.
\item Permutation symmetry of the Schwinger hierarchy gives graded
commutativity of the boundary-value fields at spacelike separation.
\item Continuing the Wightman distributions back to ordered imaginary time
recovers the original \(S_n\) away from the partial diagonals.
\end{enumerate}
\end{proposition}

\begin{proof}
For (a), the tempered tube bounds are exactly the hypotheses of the standard
boundary-value theorem for holomorphic functions of polynomial growth in a
tube.  Explicitly, for a Schwartz test \(f\), the pairings
\(\langle \mathcal W_n(\cdot+\ii y),f\rangle\) are Cauchy as
\(y\to0\) inside every smaller cone because Cauchy's integral formula bounds
the change in \(y\) by the same finite set of Schwartz seminorms.  The
resulting limit is continuous in \(f\), hence a tempered distribution.

For (b), take two adjacent insertions whose Lorentzian test supports are
spacelike separated.  The two Euclidean orderings determine two holomorphic
functions, one for each order.  By permutation symmetry they agree on the
common Euclidean set, with the graded Koszul sign.  The Jost-domain
connectivity in Definition~\ref{def:os-lorentzian-boundary-value-package}
places the two functions on opposite sides of a common real edge.  The
edge-of-the-wedge theorem identifies their analytic continuations in the
envelope, and therefore identifies their Lorentzian boundary values.  This
boundary-value identity is precisely graded commutativity of the smeared
fields for spacelike separated supports.

For (c), both the original Schwinger function on the ordered Euclidean domain
and the function obtained by first taking the Lorentzian boundary value and
then substituting ordered imaginary times are restrictions of the same
holomorphic function \(\mathcal W_n\) on a connected tube domain.  The
identity theorem for holomorphic functions gives equality away from the
partial diagonals.  Contact-term extensions on the diagonals are additional
Euclidean data and are not recovered from this off-diagonal continuation
alone.
\end{proof}

\section{Reconstruction Theorem}

\begin{theorem}[Osterwalder--Schrader reconstruction]
\label{thm:os-reconstruction}
Let \(\{S_n\}_{n\ge0}\) be Euclidean distributions satisfying Euclidean
covariance, permutation symmetry, reflection positivity, Hermiticity,
normalization, the OS-II positive-time semigroup conditions
\eqref{eq:os-time-semigroup-law}--\eqref{eq:os-time-strong-continuity}, and a
linear-growth estimate such as \eqref{eq:os-linear-growth-condition} strong
enough to prove the Lorentzian boundary-value package of
Definition~\ref{def:os-lorentzian-boundary-value-package}.  Then there exists
Wightman field data
\[
  (\Hilb_{\rm OS},\vac_{\rm OS},U,\mathcal D_{\rm OS},
  \{\widehat\Phi_\alpha\}_{\alpha\in I})
\]
whose Schwinger functions are \(\{S_n\}\) at noncoincident Euclidean points.
The reconstructed Wightman functions are obtained as boundary values of
analytic continuations of the \(S_n\).
The symbol \(\vac_{\rm OS}\) denotes the vacuum vector produced by this
reconstruction; it is a subscripted instance of the global vacuum notation
\(\vac\).
\end{theorem}

\begin{proof}
Let \(\mathcal E_+\) be the vector space of finite labelled Euclidean
insertions whose time coordinates are strictly positive.  Reflection
positivity makes
\[
  \langle F,G\rangle_{\rm OS}:=S_{m+n}(\Theta F,G)
\]
a positive semidefinite Hermitian form on \(\mathcal E_+\).  Quotienting by
the null space
\[
  \mathcal N_{\rm OS}
  =
  \{F\in\mathcal E_+:\langle F,F\rangle_{\rm OS}=0\}
\]
and completing gives the Hilbert space
\(\Hilb_{\rm OS}=\overline{\mathcal E_+/\mathcal N_{\rm OS}}\).  The class of
the empty insertion is the vector \(\vac_{\rm OS}\), and finite positive-time
insertions are dense by construction.

For \(s\ge0\), Euclidean time translation by \(s\) maps
\(\mathcal E_+\) to itself.  The OS-II semigroup assumptions
\eqref{eq:os-time-semigroup-law}--\eqref{eq:os-time-strong-continuity}
therefore descend to a strongly continuous contraction semigroup
\(T_E(s)\) on \(\Hilb_{\rm OS}\).  The symmetry
\eqref{eq:os-time-symmetry}, as verified in the positive-time quotient above,
makes each \(T_E(s)\) self-adjoint.  The self-adjoint
contraction-semigroup theorem gives a nonnegative self-adjoint generator
\[
  H,\qquad T_E(s)=e^{-sH},\qquad \sigma(H)\subset[0,\infty).
\]
Stone's theorem then gives Lorentzian time translations
\(U(t,0)=e^{itH}\).  Spatial translations and rotations preserve the
positive-time half-space after a common positive time shift; by Euclidean
covariance and the same null-space argument they produce unitary
representations on the reconstructed Hilbert space.  The OS-II
boundary-value package, supplied by the declared growth estimate, gives the
tempered Lorentzian boundary values and the analytic continuation of
Euclidean rotations mixing time and space to real Lorentzian boosts.  This
gives a unitary representation \(U\) of the connected Poincare group whose
joint translation spectrum is contained in the closed forward light cone.

For a labelled test function \(f\) supported at positive Euclidean time,
define the field insertion on the dense domain
\(\mathcal D_{\rm OS}\subset\Hilb_{\rm OS}\) represented by finite
positive-time insertions by concatenation:
\[
  \widehat\Phi_\alpha(f)\,[F]
  :=
  [(\alpha,f),F],
\]
after translating the supports, if necessary, so that the concatenated
configuration remains strictly time ordered before taking the boundary value.
The continuity estimates in the OS hypotheses make the resulting matrix
elements distributions in \(f\).  Their Lorentzian boundary values are the
Wightman distributions.  Proposition~\ref{prop:os-boundary-package-consequences}
then gives local graded commutativity for spacelike separated Lorentzian test
supports.
Hermiticity of the Schwinger hierarchy gives the adjoint relation on
\(\mathcal D_{\rm OS}\), and Euclidean covariance gives the Poincare
covariance of the reconstructed fields.

In the Euclidean time-ordered domain with noncoincident insertion points,
continuing the Wightman distributions back to imaginary time recovers the
original \(S_n\) by Proposition~\ref{prop:os-boundary-package-consequences}.
The assertion is restricted to noncoincident Euclidean points because contact
extensions are part of the declared Schwinger hierarchy rather than a
consequence of analytic continuation away from the diagonals.  The uniqueness
statement is the GNS-type argument in
Proposition~\ref{prop:os-reconstruction-uniqueness}.
\end{proof}

\begin{proposition}[Uniqueness of OS reconstruction]
\label{prop:os-reconstruction-uniqueness}
Two OS reconstructions from the same Schwinger hierarchy are unitarily
equivalent on the reconstructed Wightman field presentation.
\end{proposition}

\begin{proof}
Let
\[
  (\Hilb_1,\vac_1,U_1,\mathcal D_1,\widehat\Phi^{(1)}_\alpha),
  \qquad
  (\Hilb_2,\vac_2,U_2,\mathcal D_2,\widehat\Phi^{(2)}_\alpha)
\]
be two reconstructions from the same hierarchy.  Both start from the same
positive-time vector space \(\mathcal E_+\) and the same OS sesquilinear
form, hence from the same null space \(\mathcal N\).  On the dense quotient
define
\[
  V_0[F]_1=[F]_2 .
\]
This is well defined because \([F]_1=0\) iff
\(\langle F,F\rangle_{\rm OS}=0\) iff \([F]_2=0\).  It is isometric because
both inner products are the same OS form:
\[
  \langle V_0[F]_1,V_0[G]_1\rangle_{\Hilb_2}
  =
  \langle F,G\rangle_{\rm OS}
  =
  \langle [F]_1,[G]_1\rangle_{\Hilb_1}.
\]
Since the quotient is dense in both completions, \(V_0\) extends uniquely to a
unitary \(V:\Hilb_1\to\Hilb_2\).  The empty insertion is mapped to the empty
insertion, so \(V\vac_1=\vac_2\).  Euclidean translations, rotations, and the
continued Poincare transformations are all defined on the dense quotient by
the same operations on test sequences, hence
\[
  VU_1(g)=U_2(g)V
\]
on a dense domain and therefore on the Hilbert spaces.  The same argument
applies to field insertions:
\[
  V\widehat\Phi^{(1)}_\alpha(f)[F]_1
  =
  V[(\alpha,f),F]_1
  =
  [(\alpha,f),F]_2
  =
  \widehat\Phi^{(2)}_\alpha(f)V[F]_1,
\]
with equality on the common invariant domain.  Thus the two reconstructions
are unitarily equivalent as Wightman field presentations.
\end{proof}

\begin{corollary}[path-integral Schwinger limits as Wightman data]
\label{cor:path-integral-schwinger-limit-to-wightman}
Let \(\{\mathcal E_\Lambda,\langle\,\cdot\,\rangle_\Lambda\}\) be a regulated
Euclidean path-integral presentation, where \(\Lambda\) denotes the regulator
and \(\langle\,\cdot\,\rangle_\Lambda\) is the regulated expectation
functional on a declared algebra of Euclidean field insertions.  Assume that
for every \(n\) and every labelled test insertion \(f_n\) in the chosen test
class the regulated moments
\[
  S_{n,\Lambda}(f_n)
  :=
  \langle f_n(\Phi_E,\ldots,\Phi_E)\rangle_\Lambda
\]
converge to a distributional limit \(S_n(f_n)\), and that the limiting
hierarchy \(\{S_n\}_{n\ge0}\) satisfies the OS-admissibility hypotheses of
Definition~\ref{def:os-admissible-euclidean-data}.  Then the continuum path-integral
presentation determines a Wightman field presentation by applying
Theorem~\ref{thm:os-reconstruction} to the limiting Schwinger hierarchy.
If the limiting hierarchy also satisfies the Euclidean cluster condition in
Corollary~\ref{cor:os-reconstruction-vacuum-sector}, the reconstructed
zero-momentum subspace is \(\mathbb C\vac_{\rm OS}\).
\end{corollary}

\begin{proof}
The convergence assumption defines a family
\(\{S_n\}_{n\ge0}\) of continuous linear functionals on the declared
test-function spaces; \(S_0=1\) is the normalization of the regulated
expectation functional after taking the limit.  The OS-admissibility
assumption supplies Euclidean covariance, permutation symmetry, Hermiticity,
reflection positivity, and the OS-II semigroup and linear-growth estimates used in
Theorem~\ref{thm:os-reconstruction}.  Applying that theorem constructs
\[
  (\Hilb_{\rm OS},\vac_{\rm OS},U,\mathcal D_{\rm OS},
  \{\widehat\Phi_\alpha\})
\]
whose Wightman boundary values have the limiting Schwinger hierarchy as their
Euclidean continuation at noncoincident points.  The last sentence is exactly
Corollary~\ref{cor:os-reconstruction-vacuum-sector}.
\end{proof}

\begin{remark}[Constructive status of OS inputs]
\label{rem:os-constructive-status-inputs}
OS reconstruction is a bridge theorem.  It turns a Schwinger hierarchy with
reflection positivity, Euclidean covariance, permutation symmetry,
Hermiticity, clustering when desired, and the OS-II analytic bounds into a
Lorentzian Wightman presentation.  It does not by itself construct the
Schwinger hierarchy.  Completed constructions verifying such inputs are
model-dependent and concentrated in low-dimensional regimes, including
\(P(\phi)_2\), massive \(\phi^4_3\), selected two-dimensional
scalar--fermion models, and two-dimensional Yang--Mills in holonomy or
heat-kernel formulations.  Standard reflection-positive
Ising/\(\lambda\phi^4\) routes in \(D\ge4\) are constrained by triviality
theorems, while four-dimensional Yang--Mills, QCD-like theories, and the
full Standard Model remain open as complete constructive continuum
theorems.  Chapter~\ref{ch:constructive-status-routes} gives the status table
and route-by-route discussion.
\end{remark}

\begin{corollary}[OS reconstruction with and without clustering]
\label{cor:os-reconstruction-vacuum-sector}
Let \(P_0^{\rm OS}\) be the joint spectral projection of the reconstructed
Lorentzian translation representation onto zero momentum.  The reconstruction
theorem by itself produces an invariant cyclic vector \(\vac_{\rm OS}\) and a
positive-energy representation, but it leaves the dimension of
\(P_0^{\rm OS}\Hilb_{\rm OS}\) to the Schwinger hierarchy.  Thus, in the
absence of a cluster condition, the reconstructed representation may contain
several translation-invariant vacuum vectors.

If the Euclidean hierarchy clusters, then for finite polynomial Euclidean
field insertions \(A_E\) and \(B_E\) separated by a large spatial Euclidean
translation,
\[
  \lim_{\lambda\to+\infty}
  \langle A_E\,B_E(\lambda a)\rangle_E
  =
  \langle A_E\rangle_E\,\langle B_E\rangle_E,
  \qquad a\in\mathbb R^{D-1},\ a\ne0,
\]
then its Lorentzian reconstruction satisfies the Wightman weak cluster
property.  Consequently
\[
  P_0^{\rm OS}\Hilb_{\rm OS}=\mathbb C\vac_{\rm OS}
\]
by Theorem~\ref{thm:wightman-cluster-vacuum-uniqueness}.  Conversely, a
one-dimensional reconstructed zero-momentum subspace gives Wightman
clustering, and its restriction back to noncoincident Euclidean points gives
the corresponding Euclidean cluster property.
\end{corollary}

When \(P_0^{\rm OS}\Hilb_{\rm OS}\) has dimension larger than one, the
Schwinger functional is a translation-invariant state that is not extremal. Its
central decomposition is a direct integral, or in finite mixtures a convex
sum, of pure vacuum sectors.  OS positivity and covariance still reconstruct a
valid Wightman presentation, while clustering is the additional condition that
selects a single pure infinite-volume phase.

The theorem has three distinct pieces.  First, reflection positivity constructs
the Hilbert space, vacuum vector, and positive Hamiltonian by the positive-time
quotient.  Second, Euclidean covariance and the OS regularity assumptions
construct the Lorentzian Poincare representation and its spectrum condition.
Third, growth and analyticity assumptions control the distributional boundary
values and produce Wightman functions in the intended test-function class.  In
the tempered case, this target class is \(\mathcal S'(M^n)\).  In larger
analytic frameworks, such as hyperfunction versions of reconstruction, the
target distribution class is correspondingly larger and the growth hypotheses
are changed.

The uniqueness statement is internal to the specified Euclidean hierarchy,
label space, reflection map, and test-function class.  Two reconstructions of
the same OS-admissible data are related by a unitary equivalence preserving the
vacuum, the reconstructed Poincare representation, and the smeared fields.
Equivalence of two different Euclidean presentations, or of the local nets
generated after taking closures and affiliations, is an additional comparison
problem.

\section{Two-Point Spectral Positivity}

For a scalar field, OS positivity gives a direct Euclidean form of spectral
positivity.  Let \(f(\mathbf x)\) be a spatial test function and let
\(\Phi_E(\tau,f)\) denote the reconstructed Euclidean field smeared at
positive time \(\tau\).  The notation \(\Phi_E(0,f)\) below denotes the
boundary value from positive Euclidean time when this domain limit exists.
For \(\tau>0\), time translation in the reconstructed Hilbert space gives
\[
  \langle \vac_{\rm OS},
  \Phi_E(\tau,f)\Phi_E(0,g)\vac_{\rm OS}\rangle_{\rm OS}
  =
  \langle \Phi_E(0,f)^\dagger\vac_{\rm OS},
  \ee^{-\tau H}\Phi_E(0,g)\vac_{\rm OS}\rangle_{\rm OS}.
\]
Applying the spectral theorem to the nonnegative operator \(H\), this is
\[
  \int_{[0,\infty)} \ee^{-\tau E}\,
  \dd\nu_{f,g}(E),
\]
where
\[
  \nu_{f,g}(\Delta)
  =
  \langle \Phi_E(0,f)^\dagger\vac_{\rm OS},
  E_H(\Delta)\Phi_E(0,g)\vac_{\rm OS}\rangle_{\rm OS}
\]
and \(E_H\) is the spectral measure of \(H\).  For \(f=g\), the measure is
positive.  Including spatial translations refines this to the joint spectral
measure of \((H,\mathbf P)\).  Lorentz covariance and the spectrum condition
then repackage the scalar two-point measure by invariant mass, giving the
K\"all\'en--Lehmann representation used in Chapter~\ref{ch:kallen-lehmann}.

This calculation identifies the precise bridge: reflection positivity builds
the Hilbert-space inner product, the Euclidean time semigroup gives the
nonnegative Hamiltonian, and the spectral theorem supplies the positive
energy-momentum measure.

\section{Euclidean Measures and Generating Functionals}

Suppose a Euclidean theory is given by a probability measure \(\mu\) on a
specified measurable space of scalar distributions \(\varphi\), with
generating functional
\[
  Z_E[J]=\int \exp\left(\int_E J(x)\varphi(x)\,\dd^Dx\right)\,\dd\mu(\varphi)
\]
for real test functions \(J\) in a neighborhood of the origin.  The Schwinger
functions are moments:
\[
  S_n(x_1,\ldots,x_n)
  =
  \int \varphi(x_1)\cdots\varphi(x_n)\,\dd\mu(\varphi).
\]
Reflection positivity is then a property of the measure together with the
chosen time reflection.  It is stable under lattice approximations only when
the regulator action, observable insertions, and continuum limit preserve the
positive-time quadratic form.

\begin{remark}
Ordinary positivity of the probability measure \(\mu\) means
\(\int |F(\varphi)|^2\,\dd\mu(\varphi)\ge0\) for measurable functions \(F\).
OS positivity is the reflected statement
\(\int \overline{F(\theta\varphi)}F(\varphi)\,\dd\mu(\varphi)\ge0\) for
positive-time cylinder functions \(F\), with the label adjunction included.
It is this reflected positivity that becomes the Lorentzian Hilbert-space
inner product.
\end{remark}

\section{Reflection Positivity for Generating Functionals}

When the Euclidean moments are encoded by \(Z_E[J]\), reflection positivity
has an equivalent source form.  Let \(J_r\) be finitely many sources supported
in positive Euclidean time and let \(c_r\in\mathbb C\).  Define the reflected
source by
\[
  (\Theta J_r)(x)=\overline{J_r(\theta x)}
\]
with the field adjunction included when the source carries labels.  The
generating functional is reflection positive when
\begin{equation}
  \sum_{r,s}\overline c_r c_s\,
  Z_E[\Theta J_r+J_s]\ge0 .
  \label{eq:os-generating-functional-positivity}
\end{equation}
Expanding \(Z_E\) in powers of the sources recovers
\eqref{eq:os-reflection-positivity} for all Schwinger functions.  Conversely,
if the moment expansion is valid in a neighborhood of the origin, the
Schwinger-function inequalities imply
\eqref{eq:os-generating-functional-positivity}.

For a Gaussian scalar measure with covariance \(C(x-y)\), the criterion
reduces to positivity of the kernel reflected across \(\tau=0\):
\[
  \int_{\tau_x>0}\dd^Dx
  \int_{\tau_y>0}\dd^Dy\,
  \overline{f(x)}\,C(\theta x-y)\,f(y)\ge0
\]
for all positive-time test functions \(f\).  The free massive covariance
\[
  C(x-y)=
  \int\frac{\dd^Dk_E}{(2\pi)^D}
  \frac{\ee^{\ii k_E\cdot(x-y)}}{k_E^2+m^2}
\]
satisfies this condition because its Fourier representation in spatial
momentum is
\[
  C(\tau,\mathbf x)
  =
  \int\frac{\dd^{D-1}\mathbf p}{(2\pi)^{D-1}}
  \frac{\ee^{\ii\mathbf p\cdot\mathbf x}
        \ee^{-\omega_{\mathbf p}|\tau|}}
       {2\omega_{\mathbf p}},
  \qquad
  \omega_{\mathbf p}=\sqrt{\mathbf p^2+m^2},
\]
so the reflected quadratic form is an integral of absolute squares with
positive weight.  This computation is the Euclidean shadow of the positive
one-particle Hilbert space of the free scalar field.

\begin{figure}[H]
  \centering
  \resizebox{0.86\textwidth}{!}{%
  \begin{tikzpicture}[x=1cm,y=1cm,every node/.style={font=\scriptsize}]
    \tikzset{
      box/.style={draw,rounded corners=4pt,line width=0.8pt,fill=orange!7,
        minimum width=3.1cm,minimum height=0.85cm,align=center},
      arrow/.style={-{Latex[length=2mm]},line width=0.85pt}
    }
    \node[box] (measure) at (-4.1,0.9) {Euclidean measure\\\(\mu\)};
    \node[box] (moments) at (0,0.9) {moments\\\(S_n\)};
    \node[box] (os) at (4.1,0.9) {OS-positive\\Schwinger data};
    \node[box,fill=green!7] (wightman) at (0,-1.15)
      {Lorentzian Wightman\\boundary values};
    \draw[arrow] (measure)--(moments);
    \draw[arrow] (moments)--(os)
      node[midway,above=5pt,fill=white,inner sep=1.2pt] {verify axioms};
    \draw[arrow] (os)--(wightman);
    \draw[arrow] (moments)--(wightman)
      node[midway,left] {analytic continuation};
  \end{tikzpicture}
  }
  \caption{A Euclidean measure yields Lorentzian Wightman data through the
  verified OS positivity, covariance, and growth requirements.}
  \label{fig:euclidean-measure-os-data}
\end{figure}

\section{Relation To Wick Rotation}

For time-ordered Green functions in perturbation theory, Wick rotation was
used earlier as a contour and boundary-value operation.  OS reconstruction is
the nonperturbative counterpart at the level of complete Euclidean
correlation data.  It supplies the Hilbert space, Hamiltonian, vacuum, and
Lorentzian Wightman functions from Euclidean data satisfying precise
positivity and growth conditions.

The framework identifies the Euclidean information that reconstructs a
Lorentzian theory: a full compatible hierarchy of Schwinger distributions
together with covariance, reflection positivity, growth, regularity, and
normalization conditions.
